\begin{frame}[fragile]{Índices: o limite é responsabilidade do programa}
\begin{lstlisting}
int leituras[5] = {26, 28, 30, 32, 34};
int indice = 4;
if (indice >= 0 && indice < 5) {
    printf("Leitura: %d\n", leituras[indice]);
} else {
    printf("Indice invalido\n");
}
\end{lstlisting}
\begin{alertblock}{Acesso fora dos limites}
\texttt{leituras[5]} e \texttt{leituras[-1]} são inválidos. Em C, esse acesso tem comportamento indefinido: não existe garantia de erro visível, proteção automática ou resultado previsível.
\end{alertblock}
Para um vetor de capacidade $N$, o intervalo válido é $0 \leq i < N$.
\end{frame}
